\section*{3月6日：二次根式的基础应用1}


\begin{enumerate}
    \item 当$x$为怎样的实数时，下列各式在实数范围内有意义？
    
    (1) $\dfrac{3-x}{\sqrt{x-4}}+(x-6)^0$；
    ~~
    (2) $\sqrt{(2-x)^2}$；
    \item 写出代数式$\sqrt{x^2}$和$(\sqrt{x})^2$中$x$的取值范围，并化简代数式.
    \item 已知$a, b$为等腰$\triangle ABC$两边的长，且满足
    $b=5+\sqrt{4-2a}+\sqrt{a-2},$
    求$\triangle ABC$周长.
\end{enumerate}


\section*{3月8日：二次根式的基础应用2}

\begin{enumerate}
    \item 当$x$取$\_\_\_\_\_\_$时，$\sqrt{3x-5}$的值最小，此时$\sqrt{3x-5}=\_\_\_\_\_\_$.
    \item 当$x$取$\_\_\_\_\_\_$时，$-\sqrt{2x-1}$的值最大，此时$-\sqrt{2x-1}=\_\_\_\_\_\_$.
    \item 若$\sqrt{3-x}$为整数，则正整数$x$的值为$\_\_\_\_\_\_$.
    \item 若$\sqrt{24n}$为整数，则正整数$n$的最小值为$\_\_\_\_\_\_$.
\end{enumerate}


\section*{3月9日：二次根式的双重非负性plus}

\begin{enumerate}
\item 要使$\sqrt{3-x}+\dfrac{(x-1)^0}{\sqrt{x+6}}$有意义，则$x$应满足$\_\_\_\_\_\_$.
\item 要使$\dfrac{\sqrt{|x-3|-2}}{x^2-4x+3}$有意义，则$x$应满足$\_\_\_\_\_\_$.
\item 要使$\dfrac{\sqrt{3-|x-1|}}{\sqrt{|x-1|-2}}$有意义，则$x$应满足$\_\_\_\_\_\_$.
\item 要使$\sqrt{\dfrac{x-2}{3-x}}=\dfrac{\sqrt{x-2}}{\sqrt{3-x}}$成立，则$x$应满足$\_\_\_\_\_\_$.
\end{enumerate}




\section*{3月10日：二次根式的性质与绝对值的化简综合}

\begin{enumerate}
    \item 若$\sqrt{(x-1)^2}+|x-2|$的化简结果为$2x-3$，则$x$的取值范围为$\_\_\_\_\_\_$.
    \item 若$\sqrt{(1-a)^2}+\sqrt{(3-a)^2}$的值为$2$，则$a$的取值范围为$\_\_\_\_\_\_$.
    \item 若$\sqrt{x^2+6x+9}+\sqrt{x^2-2x+1}-\sqrt{x^2+4x+4}=-x+2$，则$x$的取值范围为$\_\_\_\_\_\_$.
\end{enumerate}

\newpage

\section*{3月11日：二次根式的乘除运算与最简二次根式的化简}

\begin{enumerate}


\item （1）$4\sqrt{\dfrac{1}{3}}\times(-2\sqrt{54}).$
（2）$\sqrt{\dfrac{1}{2}}\times\sqrt{2 \dfrac{1}{2}}\times\sqrt{1 \dfrac{3}{5}}.$
（3）$4\sqrt{27}\div 2\sqrt{\dfrac{3}{8}}.$

\item（1）$-\sqrt{2}\div\dfrac{3}{5}\sqrt{\dfrac{1}{3}}\div\sqrt{\dfrac{1}{4}}.$
（2）$\sqrt{\dfrac{2}{7}}\div\left(-\dfrac{2}{3}\sqrt{\dfrac{3}{5}}\right)\times\dfrac{1}{2}\sqrt{\dfrac{7}{5}}.$

\end{enumerate}





\section*{3月12日：含字母的二次根式的计算与化简}

\begin{enumerate}
    \item 计算：$\sqrt{\dfrac{1}{3}ab^2}\times\sqrt{\dfrac{15a}{b}}. (a\geq 0, b>0)$
    \item 计算：$\sqrt{\dfrac{5a^3b^2}{3c^4}}\div\sqrt{ab^2}\times\sqrt{\dfrac{2c}{4ab^3}}.$
    \item 若$\sqrt{a^3+a^2}=-a\sqrt{a+1}$，则$a$的取值范围为$\_\_\_\_\_\_$.
\end{enumerate}

\section*{3月13日：二次根式背景下两个很好的二级结论题}

\begin{enumerate}
    \item 海伦在其著作《度量》中给出了三角形面积计算公式（海伦公式）：$$S=\sqrt{p(p-a)(p-b)(p-c)}, ~p=\dfrac{a+b+c}{2}$$其中$a, b, c$分别是三角形的三边长。如果一个三角形的三边长分别为$7, 8, 9$，那么它的面积为$\_\_\_\_\_\_$. 
    \item 已知：$$\sqrt{1-\dfrac{1}{2}}=\sqrt{\dfrac{1}{2}}, \sqrt{2-\dfrac{2}{5}}=2\sqrt{\dfrac{2}{5}}, \sqrt{3-\dfrac{3}{10}}=3\sqrt{\dfrac{3}{10}},$$请用字母描述该规律. 
\end{enumerate}




\section*{3月14日：二次根式的综合运算}

\begin{enumerate}
    \item 计算：$\sqrt{24}\times\sqrt{\dfrac{1}{3}}+\sqrt{4}\times\sqrt{\dfrac{1}{8}}\times(\sqrt{2}-\sqrt{3})+\dfrac{1}{\sqrt{3}-2}.$
    \item 计算：$\sqrt{27a}-a\sqrt{\dfrac{3}{a}}+3\sqrt{\dfrac{a}{3}}+\dfrac{1}{2a}\sqrt{75a^3}.$
\end{enumerate}


\newpage


\section*{3月15日：特殊的二次根式化简}

\begin{enumerate}
    \item 化简：$\sqrt{3+2\sqrt{2}}.$
    \item 化简：$\sqrt{27-\sqrt{48}}.$
    \item 化简：$\sqrt{8+\sqrt{63}}+\sqrt{8-\sqrt{63}}.$
    \item 化简：$2\sqrt{3-2\sqrt{2}}+\sqrt{17-12\sqrt{2}}.$
\end{enumerate}




\section*{3月16日：二次根式中的待定系数法}

\begin{enumerate}
    \item 已知$m, n$是有理数，且有$$(\sqrt{5}+2)m+(3-2\sqrt{5})n+7=0,$$则$m=\_\_\_\_\_\_, n=\_\_\_\_\_\_.$
    \item 已知$a, b$是有理数，且有$$\left(\dfrac{1}{3}+\dfrac{\sqrt{3}}{2}\right)a+\left(\dfrac{1}{4}-\dfrac{\sqrt{3}}{12}\right)b-2\dfrac{1}{4}-1\dfrac{9}{20}\sqrt{3}=0,$$求$a, b$的值.
\end{enumerate}


\section*{3月17日：复杂二次根式混合运算综合（一）}

\begin{enumerate}
    \item $\left(\dfrac{2}{3}\sqrt{5}+2\right)\times\left(\dfrac{2}{3}\sqrt{5}-2\right)-(1-\sqrt{5})^2-\sqrt{4\dfrac{1}{3}}.$
    \item $(2\sqrt{2}-1)^2-\left(2\sqrt{2}+\dfrac{1}{2}\sqrt{3}\right)\times\left(2\sqrt{2}-\dfrac{1}{2}\sqrt{3}\right)-\left|5\sqrt{2}-3\sqrt{6}\right|.$
\end{enumerate}

\section*{3月18日：复杂二次根式混合运算综合（二）}

\begin{enumerate}
    \item $\left(2\sqrt{2}-\sqrt{4\dfrac{1}{2}}\right)^2+|2\sqrt{5}-4\sqrt{2}|-(1-\pi)^0.$
    \item $|5\sqrt{8}-3\sqrt{13}|-\left(3\sqrt{5}-\sqrt{9\dfrac{1}{2}}\right)^2-(3\sqrt{3}+\sqrt{2})\times(3\sqrt{3}-\sqrt{2}).$
    \item $(2\sqrt{3}+\sqrt{5})\times(\sqrt{5}-2\sqrt{3})-|3\sqrt{2}-2\sqrt{5}|-(\sqrt{5}-\sqrt{8})^2.$
\end{enumerate}


\newpage


\section*{3月19日：二次根式中的比大小}

\begin{enumerate}
    \item 比较$-3\sqrt{5}$和$-5\sqrt{3}$的大小.
    \item 比较$\sqrt{2}+\sqrt{6}$和$\sqrt{3}+\sqrt{5}$的大小.

\end{enumerate}






\section*{3月20日：裂项相消2.0}

\begin{enumerate}
    \item 已知：$$\dfrac{3}{\sqrt{5}-\sqrt{2}}=\dfrac{3(\sqrt{5}+\sqrt{2})}{(\sqrt{5}-\sqrt{2})(\sqrt{5}+\sqrt{2})}=\dfrac{3\sqrt{5}+\sqrt{2}}{3}=\sqrt{5}+\sqrt{2}.$$
    （1）化简$\dfrac{1}{\sqrt{3}+\sqrt{2}}.$

    （2）计算：$\dfrac{1}{1+\sqrt{2}}+\dfrac{1}{\sqrt{2}+\sqrt{3}}+\dfrac{1}{\sqrt{3}+\sqrt{4}}+\cdots+\dfrac{1}{\sqrt{99}+\sqrt{100}}.$

    （3）计算：$\dfrac{1}{1-\sqrt{2}}+\dfrac{1}{\sqrt{2}-\sqrt{3}}+\dfrac{1}{\sqrt{3}-\sqrt{4}}+\cdots+\dfrac{1}{\sqrt{99}-\sqrt{100}}.$
\end{enumerate}



